\section{Joint crossings and prescribed regions}
\label{sec:critical-tests}

Section~\ref{sec:critical} identifies the oriented critical loops along
each deterministic good sequence.  We now prove convergence of the
observables needed for the near-critical argument: finite collections
of polygon crossings after several separated squares have been forced
black or white.  First we recover the literal crossings confined to
each polygon from the oriented loops.  We then combine continuity under
small motions of the prescribed regions with finite witnesses made of
ordinary polygon crossings.  Proposition~\ref{prop:critical-tests}
collects the resulting common joint law, including its conditional
version for Poisson geometry.

\subsection{Reading joint crossing tests from the loops}
\label{subsec:critical-readout}

An interface may enter and leave a quad repeatedly.  We therefore
determine the quad crossing from the boundary pairing of interface
portions inside it.

Write the sides of a quad in boundary order as $S_1,S_2,S_3,S_4$, with
$S_1,S_3$ the black target sides. At a finite mesh, draw the interfaces along
cell edges, retaining their orientations and hence the color on each bank.
In general position with respect to a polygon, their portions in its interior
are disjoint crosscuts and closed curves. The following elementary rule reads
the crossing from these portions.

\begin{lemma}[Boundary-pairing rule]\label{lem:boundary-pairing}
Suppose finitely many disjoint interfaces, with their two bank colors, cut a
Jordan quad, avoid its marks, and cross its boundary whenever they meet it.
If an interface crosscut joins $S_1$ to $S_3$, the black bank supplies a black
crossing. If one joins $S_2$ to $S_4$, the white bank supplies a complementary
white crossing. If neither type occurs, there is a complementary region
incident to all four sides; its color determines the crossing. Whenever a
crosscut joins the two sides incident to a mark, the orientation of the
outermost such corner cut determines this central color.
\end{lemma}

\begin{proof}
A crosscut between opposite sides separates the other pair. Paths running
arbitrarily close to its banks give the asserted color connections; they can
be taken with clearance on their respective sides. The two types of opposite
crosscut cannot coexist, because they would intersect.

In their absence every crosscut has both ends on one side or on consecutive
sides. On the boundary circle the endpoint pairing is noncrossing. A same-side
pair cuts off a pocket containing no mark. A pair on consecutive sides cuts
off the corresponding corner. Remove the pockets and, for each corner that
is cut off, its outermost corner pocket. The remaining region is connected
and incident to each of the four side intervals between the removed pockets.
This can also be proved by induction: remove an innermost paired interval
first, then restore it as a pocket along one boundary interval of the remaining
region. Interior closed curves only remove interior pockets. Consequently the
region left between the corner pockets meets all four sides. A path through
it joins either pair of sides, and its color decides which crossing occurs.
The bank of an outermost corner cut facing this region has that color.
All paths used here lie in the quad; no identification of different boundary
points or external connection is made.
\end{proof}

We explain why this finite rule is stable for the limiting loops. First put
the polygon in general position by a translation chosen from an arbitrarily
small open disk. For almost every such translation, almost surely the limit
loops avoid the vertices and marks, have no double point or mutual contact on
a side, and have no tangency there that is a local extremum of the coordinate
normal to the side. Here are the estimates and deterministic facts behind
these assertions. A prescribed point is avoided by the one-arm estimate.
A double visit or a contact of two loops on a line segment produces four
traversals in a ball centered on that segment. Covering the segment by
$O(r^{-1})$ balls and using Lemma~\ref{lem:color-arms} gives an upper bound
$O(r^{c})$, for some $c>0$, when the traversals continue to a fixed positive
distance. Let that distance run through a countable sequence and then let
$r\downarrow0$. Finally, for any continuous real function, the set of its
local-maximum and local-minimum values is countable: each such value is the
maximum or minimum on some interval with rational endpoints. Apply this to
the normal coordinate along each of the countably many loops. Almost every
translation avoids all those levels. These facts hold simultaneously for a
finite collection of polygons. General position at finite meshes is obtained
by excluding a further null set of translations; at a Voronoi vertex the
incident cells are pairwise side-adjacent, so this exclusion introduces no
new connectivity convention.

There are also corner cuts at positive, possibly random, distances from each
mark. To see this, choose many disjoint buffered annuli around the mark, all
small enough that the two incident polygon sides are the only boundary pieces
present. In each annulus, fixed crossing chains give a black circuit and a
white circuit in separated subannuli with a probability bounded below.
The colors in these buffered regions are independent on a deterministic
mesh, once the mesh is sufficiently small. Thus the probability that none
of $k$ annuli supplies the circuits is at most $(1-c)^k$. The two circuits
force an intervening closed wall enclosing the mark. Following this wall
through the polygonal sector gives a cut between the incident sides.
Take the mesh limit first and then $k\to\infty$. This proves the claim also
when the mark lies in the interior of a straight side.

Fix one such cut at each mark, and trim its two reaches on the incident sides
a positive distance away from that mark. Any outermost corner cut beyond
these reaches has positive diameter: the two compact side portions in
question are disjoint. Opposite-side excursions also have positive diameter
after the mark neighborhoods have been removed. At any fixed positive reach
there are only finitely many relevant excursions, by curve continuity and
the finite macroscopic matching in Proposition~\ref{prop:loops}. Their orders
and bank colors persist under convergence. Indeed their compact interior
portions stay in the open quad. At an endpoint on a side, the loop must
approach the line from both sides in every sufficiently small parameter
neighborhood, unless that line is a local-extremum level. Therefore nearby
curves recover the endpoint crossing. Conversely, a sequence of excursions
of positive reach has a limiting excursion: it cannot collapse at a mark,
and it cannot acquire an intermediate boundary contact without either a
double contact on the side or a forbidden one-sided tangency. Reach cutoffs
can be chosen to avoid equality. The boundary-pairing rule now gives the same
crossing decision on matched configurations. Unmatched loops of vanishing
diameter cannot change it, because the positive corner reaches have already
fixed the central color in the case with no opposite-side excursion.

For a specified, unperturbed polygon use the inner and outer Jordan sandwich
tests from Lemma~\ref{lem:good-sequences}. They are chosen so that, on the same
colors, the lower crossing implies the specified crossing, which implies
the upper crossing. The tests may be put in the preceding general position.
The difference between their limiting probabilities tends to zero with the
sandwich width, by continuity of Cardy's formula under convergence of the
marked Jordan boundaries. A union bound applies to any finite family. We
have therefore proved joint convergence of the exact closed-cell crossing
indicators, along every deterministic good sequence, to a common critical
law. This argument also proves convergence for polygonal site-path tests
with the terminal conventions of Section~\ref{sec:critical}: the positive
sandwich margins absorb the conversion between site paths and cell paths.

\subsection{Continuity under changes of prescribed regions}
\label{subsec:forced-continuity}

Fix a finite family of closed squares
\[
 H_j=\{x\in\RR^2:\|x-z_j\|_\infty\le h_j\},\qquad h_j>0,
 \quad 1\le j\le J.
\]
Assume that the squares are mutually disjoint at positive distance and that
each is at positive distance from every polygon boundary under consideration.
For each crossing test choose, separately, whether to force all cells meeting
$H_j$ black, to force them white, or to leave them unchanged. A square outside
a particular quad is irrelevant once the mesh is small. Different tests may
make different choices while using the same original fair colors. If a cell meets two differently prescribed squares, assign it the color of
the square with smaller index; this convention is eventually irrelevant on
a good sequence, since the squares have positive separation. Throughout
this subsection the number of squares, their positive radii, and their
separations are fixed before the mesh tends to zero.

Call an arm \emph{free} when it uses only sites whose original fair colors
have not been prescribed in the test. We will use the mesh-first bounds
\begin{equation}\label{eq:forced-arm-bounds}
\begin{aligned}
 \limsup_{e\to0}\PP(\hbox{one free arm from }r\hbox{ to }R)
       &\le C(r/R)^\gamma,\\
 \limsup_{e\to0}\PP(\hbox{three free half-plane arms})
       &\le C(r/R)^\beta,
\end{aligned}
\end{equation}
where $\gamma>0$ and $\beta>1$; the second bound allows any color order and
any fixed side direction. These are consequences of
Lemmas~\ref{lem:arm-inputs} and~\ref{lem:boundary-arms}. Shifts of the bounding
line and constant factors in the radii are allowed. All probabilities in the
next two lemmas are color probabilities on a deterministic good sequence.

\begin{lemma}[Continuity of prescribed-region tests]
\label{lem:forced-continuity}
For any fixed finite family of polygon crossing tests with the square
prescriptions above, the probability that a test changes under a sufficiently
small polygonal displacement of its sides, marks, or prescribed squares tends
to zero, uniformly in the mesh-first limit. The displacements may change the
centers and positive side lengths of the squares, provided the stated
separations persist. The same conclusion holds simultaneously for any fixed
finite collection of prescription choices on the same colors.
\end{lemma}

\begin{proof}
We give both the geometric interpolation and the arm test for a single local
change. Choose $u>0$ smaller than the separations of all unrelated sides and
squares, and put balls of radius $Cu$ around the polygon vertices, marks,
and square corners. Outside these balls divide each straight edge into
$O(s^{-1})$ intervals of length comparable to $s$, where $s\ll u$. For a
displacement much smaller than $s$, perform the motion one interval at a time,
using a piecewise linear cutoff supported on a slightly longer interval.
Leave the portions within distance $2u$ of the vertices for the corner
moves, and taper the edge moves between distances $2u$ and $3u$; motions in
balls of radius $4u$ then complete the interpolation. There are boundedly
many corner moves and $O(s^{-1})$ edge moves. A polygonal collar identifies
the vicinity of each edge with a straight strip. Interpolating its sufficiently
small piecewise linear ambient displacements preserves the Jordan property;
it also gives the familiar quad sandwiches by shifting target and non-target
sides in their respective directions. Squares are treated by the same
construction. Each intermediate prescribed region remains a connected disk,
and the part common to the two regions in a local comparison is connected.
This last property can be ensured by taking the motions to be graphs in their
collars and retaining the unchanged inner square. All constructions are
fixed before $e\to0$.

Consider an edge move supported in a ball of radius $Cs$. If it changes a
crossing, take a black crossing from the configuration where the crossing
occurs and an open white complementary crossing from the other configuration.
Outside the changed neighborhood their paths run in the common tiling and
coloring. For a quad-side move, their portions going to the other three sides
give three disjoint free arms from scale $Cs$ to scale $cu$. They lie in a
half-plane translated by $O(s)$ from the original supporting line. Trimming
at first and last hits of auxiliary curves preserves the order of the arms.
The paths stay in the common region of the two domains; using the enlarged
half-plane therefore does not require the intermediate boundary itself to
be straight.

At an edge of a black prescribed square, the two arms of the white
complementary path are free. At least one of the two black branches towards
the target sides is also free until scale $cu$. Otherwise both target sides
would connect, before entering the changed neighborhood, to the common
black prescribed region. That region is connected also in the comparison
without a crossing, so it would already join the target sides there, a
contradiction. The white branches cannot connect to one another in the common
part so as to avoid the changed region, since they would then block the black
crossing in the other comparison. After proper trimming we thus have three
disjoint free arms in the exterior half-plane of the square edge. A white
prescribed square is handled with the colors interchanged. The open white
connections through a white prescribed region exist with clearance, because
its forced cells form a connected union by side adjacency.

Here and below the ``cells meeting'' convention is compatible with that
connectivity. Follow a path inside the connected prescribed region through
the cells it meets. Consecutive cells share a side, or meet at a tiling vertex
where the three cells are pairwise side-adjacent. The resulting side-adjacency
chain connects all cells meeting the region. A path in those cells can be
moved into cell interiors and relative interiors of shared sides. On an
unchanged straight boundary piece, convexity of each cell and the shortcut
argument for crossing and complementary crossing remove possible reuse of a
cell by two candidate arms. The finitely many bends introduced by a fixed
interpolation require one further observation. A shared cell there, despite
the arms continuing to a fixed outer scale, would give a free arm from a
vanishing neighborhood of that bend to that scale. Its probability tends to
zero by the one-arm bound. Thus the graph-arm extraction is valid with
vanishing error, including for the exact closed-cell crossing convention.

At a corner move, a changed crossing requires at least one free arm from
scale $Cu$ to a fixed positive scale. At a quad vertex or mark this is a
branch of the crossing towards a target arc away from that vertex. At a
prescribed square it is a branch of the color opposite to that square;
this branch cannot use the prescribed region. Since all other prescribed
regions and unrelated boundary pieces are at fixed positive distances,
the arm is free throughout the annulus used for this estimate.

A union bound and \eqref{eq:forced-arm-bounds} give, for displacement small
enough in terms of $u,s$,
\begin{equation}\label{eq:forced-continuity-error}
 \limsup_{e\to0}\PP(\hbox{some compared indicator differs})
 \le C u^\gamma+C s^{-1}(s/u)^\beta.
\end{equation}
Constants may depend on the fixed tests and their separations. The vanishing
errors at the finitely many interpolation vertices have already disappeared
in the mesh limit. Given an error tolerance, first choose $u$ so that the
corner term is small, then choose $s$ so that the edge term is small, and
finally restrict the displacement. Equivalently, in the displayed bound the
limits are taken as $e\to0$, then $s\downarrow0$ at fixed $u$, and then
$u\downarrow0$. This proves the lemma.
\end{proof}

\subsection{Finite witnesses after prescribing square colors}
\label{subsec:forced-witnesses}

Continuity alone does not identify the law of the modified crossings: prescribing
a square does not make its neighboring crossings independent. We instead
show that a positive crossing, after a small favorable perturbation, has a
witness consisting of finitely many ordinary critical polygon crossings.
Those were identified in Subsection~\ref{subsec:critical-readout}.

For a black crossing, a \emph{favorable perturbation} enlarges every black
prescribed square, shrinks every white prescribed square, and makes a quad
sandwich favorable to black. The latter moves its target caps slightly towards
each other and its non-target sides outwards in polygonal collar coordinates.
A path crossing the original quad then crosses the perturbed quad after it
is clipped at the new caps. We use two nested perturbations, called the
middle and final perturbations, with positive slack between each pair.
All prescribed squares remain separated and off the quad boundary.

\begin{lemma}[Free segments of a terminal path]
\label{lem:free-segments}
Discard squares left unchanged in the test and squares outside its quad.
Construct a graph approximation to the middle quad with the remaining
middle squares removed. Give its vertices the original fair colors. Fill each square hole
by a fan from one new vertex having that square's prescribed color, and
attach uncolored terminals to the two target arcs. On the event that the
original modified configuration crosses, this graph has a black terminal
path for all sufficiently small meshes. It has such a path whose portions
between consecutive terminals or fan vertices form at most $J+1$ free paths.
Their joint laws, restricted to the crossing event, are tight in the curve
metric. Every subsequential limiting free path is simple after constant
parameter intervals are removed.
\end{lemma}

\begin{proof}
Approximate the outer boundary and hole boundaries in disjoint thin strips
by simple site cycles, in boundary order, using the graph approximation
construction of Section~\ref{sec:inputs}. Hole rims are free vertices.
The graph outside these cycles, with each hole triangulated by its fan,
is a finite planar triangulation. The terminal neighbors are exactly the
sites on the corresponding target arc; the terminals create no connection
between other boundary vertices.

Shadow an original closed-cell crossing by a site path with vanishing error.
Before it enters a black forced square, it crosses the rim of the enlarged
middle square through free black sites; crop there and replace its subsequent
passage through that square by the black fan. An original black path stays
away from every shrunken white hole. The quad collar slack allows the shadow
to be clipped at the middle target arcs and keeps it clear of the middle
non-target sides. This produces the claimed terminal path. These statements
only use the small mesh of cells on the fixed compact region and the strictly
positive slack, so hold for all sufficiently fine configurations on the
given crossing event.

Choose an outermost black terminal path on one side. More precisely, add an
external edge joining the terminals, take the union of all simple black
terminal paths together with that edge, and take the appropriate exterior
boundary cycle of this union. Each constituent path with the external edge
is a cycle, and these cycles share that edge; their union is two-connected.
Its exterior boundary therefore supplies a simple terminal path $P$.
No black path can leave $P$ into the exterior side and return to a distinct
vertex of $P$: taking its first departure and first return would produce a
terminal path farther on that side, contrary to the boundary construction.
We call this the no-bypass property. It is the same elementary outermost-path
construction used in Lemma~\ref{lem:rectangle-colors}.

A simple $P$ visits each fan vertex at most once and never visits a white
fan vertex. Remove the fan edges and terminal edges. The remaining segments
join distinct rim or terminal sets and consist entirely of free black sites.
There are at most $J+1$ of them. Many traversals of a fixed annulus by one
segment yield equally many disjoint free black arms, after proper trimming.
The one-color disjoint-occurrence bound and the one-arm bound control these
numbers. The successive-excursion compactness argument in
Lemma~\ref{lem:exploration-tightness} consequently applies to these segments,
also at their endpoints. Notice that the segments may be selected using all
the colors; the bound is on the existence of the disjoint arms and does not
assume freshness conditional on the selected path.

We give the additional argument that rules out a self-touch in a limit.
Suppose first that the contact lies in the bulk, at positive distance from
all rims. Two visits with a nonconstant intervening excursion produce four
proper traversals of a sufficiently small annulus, with a fixed positive
outer radius. Between the traversals, on the side forbidden by the
no-bypass property, there are at least two disjoint radial gaps. Indeed, in cyclic order the
four directed traversals have distinct sectors on their forbidden sides;
at most two traversals can border the same such sector. Each selected
gap contains a white arm. To justify the last assertion without assuming
smooth graph cuts, choose two intermediate graph cycles, with wide radial
margins. Within the gap each cycle has a transverse vertex path between its
two path banks. One way to see this is to label just the vertices of that
cycle black and apply Hex in the gap: an avoiding radial white path would
cross the separating cycle, which is impossible. Use an innermost such cut
at the inner cycle, and then one at the outer cycle. They enclose a
quadrilateral whose side banks are portions of $P$. A black crossing between
these banks would leave into the forbidden side and return to $P$; trim at
its first return to obtain a bypass. Hex therefore gives a white crossing
between the transverse cuts. With the cuts chosen away from the trimmed
ends, the departure points on $P$ are internal and both adjacent path edges
are still present, as required for the no-bypass argument. Trim the white
crossing to the smaller annulus. Doing this in the two disjoint gaps yields
four black and two white site-disjoint arms. All vertices involved are free;
no fan edge or rim has entered this annulus.

By Lemma~\ref{lem:color-arms}, the probability of such a six-arm event in any
of $O(r^{-2})$ balls of radius $r$ is $O(r^{c})$, when its outer radius is
fixed. It tends to zero. At a point of a straight rim or outer boundary,
a self-touch gives at least three free traversals in the corresponding
half-plane. There are three even if one of the two visits is an endpoint
of the segment. Covering that side by $O(r^{-1})$ balls and applying
Lemma~\ref{lem:boundary-arms} again gives a vanishing bound. At the finitely
many polygon corners, a single free arm reaching from vanishing scale to a
fixed scale already has probability tending to zero. The endpoints of a
segment lie in distinct, positively separated rim or terminal sets, so they
cannot coincide with each other. These exclusions at countably many outer
scales rule out every nontrivial repeated point of a limiting segment.
\end{proof}

The need for a middle perturbation is now visible; see
Figure~\ref{fig:forced-tube}. A free segment can touch a
middle rim or boundary. The final perturbation puts its black attachment
points inside an open black square and separates it from every final white
square. Thus the segment can be surrounded by a genuine open tube without
having to reconstruct any limiting connectivity at a contact point.

\begin{figure}[t]
\centering
\begin{tikzpicture}[scale=.88,every node/.style={font=\small}]
 \fill[black!14] (.35,.35) rectangle (2.25,2.25);
 \draw[thick] (.35,.35) rectangle (2.25,2.25);
 \draw[dashed] (.65,.65) rectangle (1.95,1.95);
 \node at (1.3,2.65) {enlarge black};
 \node at (1.3,1.45) {$H_B^{\rm mid}$};
 \draw[dashed] (4.25,.8) rectangle (5.75,2.3);
 \fill[pattern=north east lines,pattern color=black!35]
       (4.6,1.15) rectangle (5.4,1.95);
 \draw[thick] (4.6,1.15) rectangle (5.4,1.95);
 \node at (5,2.65) {shrink white};
 \node[fill=white,inner sep=1pt] at (5,1.55) {$H_W^{+}$};
 \draw[blue!22,line width=12pt,line cap=round]
       (1.95,1) .. controls (2.7,.6) and (3.35,.42) .. (4,.4)
       .. controls (5,.25) and (6.2,.4) .. (7.35,1.2);
 \draw[blue!65!black,thick]
       (1.95,1) .. controls (2.7,.6) and (3.35,.42) .. (4,.4)
       .. controls (5,.25) and (6.2,.4) .. (7.35,1.2);
 \fill (1.95,1) circle (2pt);
 \draw[thick] (7.35,.6)--(7.35,1.8);
 \node[right] at (7.35,1.8) {target collar};
 \node[blue!65!black] at (4,-.5) {tube around a free limiting arc};
 \node at (1.3,.04) {$H_B^{+}$};
\end{tikzpicture}
\caption{The middle rims are dashed and the final prescribed squares have
solid boundaries. A limiting free segment starts on the middle black rim,
strictly inside the final black square. Its tube stays clear of the final
white square. Its other endpoint is shown in a target collar. Thus an
ordinary crossing of a polygon quad in the tube can be joined through the
final black region and clipped at the target cap. The figure is schematic;
no regularity of the limiting arc beyond simplicity is used.}
\label{fig:forced-tube}
\end{figure}

\begin{lemma}[Transfer by polygon witnesses]\label{lem:forced-transfer}
Let two deterministic good sequences be coupled so that their ordinary
critical polygon crossing indicators converge jointly to the same indicators
for a countable family containing the rational polygons. For any fixed
prescribed-square test, the probability of a crossing in the first sequence
and no crossing in the final favorable perturbation of the second tends to
zero. The conclusion holds for finitely many tests and for independent
color replicas.
\end{lemma}

\begin{proof}
It suffices to consider a subsequence on which the asserted mismatch has
positive limiting probability. Include its indicator, the crossing vectors,
and the at most $J+1$ segments of Lemma~\ref{lem:free-segments} in a joint
subsequential representation. Pad the segment list when necessary and
include its finite combinatorial itinerary through the hubs as a coordinate
in a finite discrete space. In the representation this itinerary stabilizes
almost surely, though its limit may be random. Tightness gives convergence
of the segment curves;
Lemma~\ref{lem:free-segments} makes each nonempty limit a simple arc.
On the represented mismatch event their ordinary critical crossing vectors
agree in the limit, coordinate by coordinate.

For each limiting segment, choose its two end neighborhoods as follows.
An end at a black hub lies strictly inside the final enlarged black square;
an end at an outer terminal lies in the collar on the appropriate side of
the final target cap. The segment stays a positive distance from the final
white squares, because these are strictly smaller than the middle white
holes. Its non-target boundary approaches have room inside the final
non-target collar. Near the two ends, cut the simple arc at points still
within the prescribed end neighborhoods. The planar Schoenflies theorem gives a rectangular neighborhood of the
intervening simple arc, which can be chosen as a thin tube around it. Choose a polygon quad in this
tube whose target end cuts are inside those neighborhoods and whose lateral
sides leave room around the arc. A sufficiently close rational polygon
approximation preserves these properties. Uniform convergence in curve
order implies that the corresponding free segment in the first model
crosses this fixed quad for all sufficiently large indices.

The witnesses can depend on the limiting configuration, but are selected
from a countable list. On the representation, convergence of a Boolean
coordinate means eventual equality of that coordinate. Thus, simultaneously
for every member of the list, a witnessed black crossing in the first
sequence is also present in the second for all sufficiently large indices.
There are only finitely many witnesses on the event under consideration.
Their crossings in the second sequence avoid its final white squares: the
tube clearance eventually exceeds every cell diameter in the fixed region,
so none of the cells in a tube can be recolored white by a distant square. If
a crossing enters an enlarged black square, that only helps it. At a hub,
the two adjoining witnesses have endpoints in the open square, so the forced
black cells connect them there with clearance. At the outer ends, join in
their prescribed order and clip the resulting path at the final target caps.
This yields a black crossing of the final test, contradicting the represented
mismatch. The tubes need not be disjoint: their only required joins occur
inside their specified black squares, and any additional intersection of
black witnesses is harmless.

For several tests apply this argument to each positive indicator; for an
indicator equal to zero use the same implication in the reverse coupling
when deriving equality of laws. For replicas keep a separate segment list
and crossing vector in each replica. The arm bounds apply separately, and
a finite union bound proves the assertion without any assumption that
the replicas have independent geometries.
\end{proof}

\subsection{The resulting critical comparison}
\label{subsec:critical-comparison}

We collect the precise conclusion needed below. A test specification consists
of a polygon quad, a replica label, and a choice in $\{*,B,W\}$ for each of a
finite family of squares; $*$ means that the original fair colors are left
unchanged. A finite collection of specifications may repeat the same quad
with different prescriptions. Within a replica all tests use the same fair
color configuration.

\begin{proposition}[Joint critical tests with prescribed squares]
\label{prop:critical-tests}
Fix finitely many polygon quads and finitely many closed squares of positive
side length, with the squares mutually separated and separated from every
quad boundary. For any finite collection of test specifications, the joint
law of their exact closed-cell crossing indicators has a common limit along
every deterministic good sequence of either model. The limit is continuous
in the centers and positive side lengths of the squares throughout the open
range where the stated separations hold, and under sufficiently small marked
polygonal deformations of the quads. The conclusion remains valid when the
centers and side lengths converge to such an admissible fixed configuration
as the mesh tends to zero.

For Poisson--Voronoi percolation the conditional law given $\eta_e$ converges
in probability, in total variation on this finite Boolean space, to the
same deterministic limiting law. With any fixed number of conditionally
independent fair-color replicas on one Poisson geometry, the limiting law
is that of independent copies of the same critical configuration, with all
the prescribed tests performed in their specified replicas. In particular,
two replicas sharing their geometry have the same joint critical limit as
two independent triangular replicas.
\end{proposition}

\begin{proof}
For unmodified tests the assertion on a deterministic good sequence was
proved in Subsection~\ref{subsec:critical-readout}. On any subsequence,
compactness of the countable Boolean crossing space permits a representation
in which these ordinary test indicators converge coordinatewise. Couple two
candidate subsequential limits through their common unmodified crossing
vector. Lemma~\ref{lem:forced-transfer} transfers every positive prescribed
test from either sequence to the other's arbitrarily small favorable
perturbation. By Lemma~\ref{lem:forced-continuity}, replacing that perturbation
by the specified test costs a probability tending to zero with its size.
For precision, let $X_e^{a,i}$ be test $i$ in sequence $a\in\{1,2\}$,
and let $X_e^{a,i,+v}$ denote a favorable perturbation of size $v>0$.
The transfer implication in both directions gives
\[
 \limsup_{e\to0}\PP(X_e^1\ne X_e^2)
 \le \sum_i\sum_{a=1}^2\limsup_{e\to0}
       \PP(X_e^{a,i,+v}\ne X_e^{a,i}).
\]
The right side tends to zero as $v\downarrow0$ by continuity.
Thus the resulting Boolean vectors agree with probability tending to one
under the coupling, so their subsequential laws coincide. This proves existence
and commonality of the limit, including different prescriptions in different
tests on the same colors.

The continuity assertion follows directly from
Lemma~\ref{lem:forced-continuity}. For mesh-dependent centers or radii,
bracket each moving square between two fixed concentric squares with radii
arbitrarily close to its limiting radius. Force the inner or outer square
according to the monotone direction of its color: enlarging black helps a
black crossing, while enlarging white hinders it. With a corresponding quad
sandwich if needed, the specified moving test lies between two fixed tests.
Continuity makes their disagreement probability vanish. This proves the
moving-parameter statement without requiring estimates at mesh-dependent
positive scales.

On a deterministic tiling independent color replicas have product laws.
Apply the preceding convergence to the finite vector in each replica; their
joint law therefore converges to the corresponding product. For Poisson
geometry, Lemma~\ref{lem:good-sequences} says that from every vanishing-mesh
sequence one can extract a subsequence along which the environments are
almost surely good deterministic sequences. On this subsequence, each
conditional atom probability for our finite Boolean vector converges almost
surely to its just identified deterministic value. There are finitely many
atoms, so the conditional total-variation distance tends to zero almost
surely. The subsequence criterion proves convergence in probability along
the original sequence. The argument applies to the product color measure
on the same fixed environment and hence retains the shared geometry in the
replica assertion. Averaging gives the corresponding annealed limits.
\end{proof}
